\section{Quantitative transfer from the pencil to Haar measure}
\label{sec:transfer}

This section compares the empirical Frobenius classes in the pencil with
Haar measure after applying the probability-measure-valued race functional.
There are two quantitative costs: the compactly supported cohomology of a
tame character sheaf and the heat trace of \(\USp(2g)\).  We keep both costs
explicit in \(g\).

\subsection{The weight-zero sheaf and a direct character estimate}

Write \(Q=s^2\) and let
\(\pi_g:\mathcal C_g\to U_g\) be the smooth proper family obtained from
\eqref{eq:pencil}.  We use geometric Frobenius throughout: thus
\(\operatorname{Fr}^{\mathrm{geom}}_{Q,t}\) acts on
\(R^1(\pi_g)_*\mathbb Q_3\) with characteristic polynomial
\[
 \det(1-T\operatorname{Fr}^{\mathrm{geom}}_{Q,t})
   =\prod_{j=1}^{2g}(1-\alpha_{j,t}T),
 \qquad \alpha_{j,t}\alpha_{j+g,t}=Q.
\tag{4.1}\label{eq:frob-convention}
\]
Thus the normalized Frobenius class has eigenvalues \(\alpha_{j,t}/s\).
This convention is essential in the following half-Tate twist.

The arithmetic quotient of \(\pi_1(U_g)\) is procyclic.  Since
\(3\nmid s\), there is a continuous constant \(3\)-adic character
\(\eta_s\), geometrically trivial, such that
\[
 \eta_s(\operatorname{Fr}^{\mathrm{arith}}_Q)=s,
 \qquad
 \eta_s(\operatorname{Fr}^{\mathrm{geom}}_Q)=s^{-1}.
\tag{4.2}\label{eq:half-tate-character}
\]
Indeed, the arithmetic quotient maps continuously to the closed procyclic
subgroup generated by \(s\) in \(\mathbb Z_3^\times\).  Put
\[
 \mathcal F_g=R^1(\pi_g)_*\mathbb Q_3\otimes\eta_s.
\]
The original cup-product pairing has Frobenius multiplier \(Q\); after the
twist its geometric-Frobenius multiplier is
\(Q\eta_s(\operatorname{Fr}^{\mathrm{geom}}_Q)^2=1\).  The character
\(\eta_s\) has weight \(-1\) and is trivial on the geometric fundamental
group.  At a degree-\(n\) closed point it multiplies every geometric-Frobenius
eigenvalue by \(s^{-n}=Q^{-n/2}\), so the resulting sheaf is pure of weight
zero.  Consequently \(\mathcal F_g\)
is lisse of rank \(2g\), has the same geometric
monodromy \(\Sp(2g)\), and its full arithmetic image lies in \(\Sp(2g)\).
Its Frobenius conjugacy classes in the compact form \(\USp(2g)\) are exactly
the normalized classes \(\Theta_{g,t}\).  The geometric monodromy and
tameness assertions used here are
\cite[Theorem~10.1.16 and Lemma~10.1.12]{KatzSarnak1999}.

Let \(\mu_{g,Q}\) be the central probability measure on \(\USp(2g)\)
obtained by putting mass \(|U_g(\F_Q)|^{-1}\) on each Frobenius conjugacy
class and averaging within that class.  Let \(m_g^{\mathrm{Haar}}\) denote
Haar probability measure.

\begin{proposition}[Quantitative character equidistribution]
\label{prop:character-equidistribution}
If \(Q>4g\), then every nontrivial irreducible character \(\chi_\rho\)
of \(\USp(2g)\) satisfies
\[
 \left|\int\chi_\rho\,d\mu_{g,Q}\right|
 \le \frac{(2g-1)\sqrt Q}{Q-2g}\dim\rho
 \le \frac{2(2g-1)\dim\rho}{\sqrt Q}.
\tag{4.3}\label{eq:character-bound}
\]
\end{proposition}

\begin{proof}
Every irreducible representation of the compact group \(\USp(2g)\) is the
restriction of the corresponding irreducible algebraic representation
\(\rho\) of \(\Sp(2g)_{\mathbb C}\).  After choosing a \(3\)-adic model,
apply \(\rho\) to \(\mathcal F_g\).  The resulting local system
\(\rho(\mathcal F_g)\) is lisse, tame, pure of weight zero, and has rank
\(\dim\rho\).  Since the geometric monodromy is the full group \(\Sp(2g)\)
and \(\rho\) is nontrivial irreducible, it has no geometric invariants.

The base \(U_g=\mathbb A^1-\{1,\ldots,2g\}\) has
\(\chi_c(U_g)=1-2g\).  Compactly supported \(H^0\) vanishes because a global
section of a lisse sheaf on the connected nonproper curve \(U_g\) cannot have
proper support unless it is zero.  Poincare duality identifies \(H_c^2\) with
the dual of the geometric invariant space, so it also vanishes.  The tame
Grothendieck--Ogg--Shafarevich formula
\cite[Theorem~9.2.5]{KatzSarnak1999} therefore gives the exact identity
\[
 \dim H_c^1\bigl(U_g\otimes\overline{\F}_Q,
                         \rho(\mathcal F_g)\bigr)
 =(2g-1)\dim\rho.
\]
The Grothendieck--Lefschetz trace formula and Deligne's weight bound
\cite{Deligne1980} now give
\[
 \left|\sum_{t\in U_g(\F_Q)}\chi_\rho(\Theta_{g,t})\right|
 \le (2g-1)\dim\rho\,\sqrt Q.
\]
Division by \(|U_g(\F_Q)|=Q-2g\) proves the first inequality in
\eqref{eq:character-bound}; the second follows from \(Q>4g\).
\end{proof}

\subsection{Lipschitz control of the quenched law}

For \(U\in\USp(2g)\), order its positive eigenangles as
\(0\le\vartheta_1\le\cdots\le\vartheta_g\le\pi\), and define
\[
 T_{g,\lambda}(U)
 =\frac12\sum_{\epsilon=0}^1
 \Law_\Phi\left(c_{g,\lambda}^{(\epsilon)}+
       \sum_{j=1}^g b_{g,\lambda}(\vartheta_j)\cos\Phi_j\right)
 \in\Ptwo.
\tag{4.4}\label{eq:T-map}
\]
This is a central map.  We equip \(\USp(2g)\) with the bi-invariant metric
whose value at the identity is
\[
 \langle X,Y\rangle=-\frac12\operatorname{Tr}(XY).
\tag{4.5}\label{eq:group-metric}
\]
On the standard maximal torus this is exactly the Euclidean angle metric.

\begin{lemma}[Metric on the conjugacy quotient]
\label{lem:conjugacy-quotient-metric}
Let \(U,V\in\USp(2g)\), and order their positive eigenangles as
\(0\le\vartheta_1\le\cdots\le\vartheta_g\le\pi\) and
\(0\le\varphi_1\le\cdots\le\varphi_g\le\pi\).  For the metric
\eqref{eq:group-metric},
\[
 \sum_{j=1}^g|\vartheta_j-\varphi_j|^2
 =d_{\mathrm{conj}}([U],[V])^2
 \le d_G(U,V)^2.
\]
\end{lemma}

\begin{proof}
For a compact connected Lie group with a bi-invariant metric, every conjugacy
class meets a fixed maximal torus, and the induced orbit metric on the
conjugacy quotient is the Euclidean torus metric modulo the affine Weyl group.
Indeed, a shortest geodesic between the two compact orbits is orthogonal to
both by the first-variation formula; bi-invariance then places its tangent in
a common maximal torus.  Concretely,
\[
 d_{\mathrm{conj}}([U],[V])^2
 =\min_{w\in W(C_g),\ k\in(2\pi\mathbb Z)^g}
      \|\vartheta-w\varphi+k\|_2^2.
\]
For vectors in the closed alcove
\(0\le\vartheta_1\le\cdots\le\vartheta_g\le\pi\), sign changes and
\(2\pi\)-translations cannot improve on pairing two nonnegative angles, and
the rearrangement inequality shows that increasing order is the minimizing
permutation.  Direct matrix
multiplication in \eqref{eq:group-metric} gives the Euclidean angle norm,
proving the equality.  The inequality follows because the orbit-distance
infimum includes the original representatives \(U,V\).
\end{proof}

\begin{lemma}[Same-phase Wasserstein coupling]
\label{lem:T-lipschitz}
For ordered positive angle vectors \(\vartheta,\varphi\),
\[
 \Wtwo\bigl(T_{g,\lambda}(\vartheta),T_{g,\lambda}(\varphi)\bigr)^2
 \le \frac12\sum_{j=1}^g
 \left|b_{g,\lambda}(\vartheta_j)
             -b_{g,\lambda}(\varphi_j)\right|^2.
\tag{4.6}\label{eq:W2-coupling}
\]
Moreover,
\[
 \operatorname{Lip}(T_{g,\lambda})
 \le \frac{\|b'_{g,\lambda}\|_\infty}{\sqrt2}
 \le \frac{\sqrt2(2g+1+\lambda)}{\lambda}.
\tag{4.7}\label{eq:T-Lip}
\]
\end{lemma}

\begin{proof}
Choose the same fair parity variable and the same phases in the two laws.
The centers cancel, while independence and
\(\mathbb E\cos^2\Phi_j=1/2\) give exactly the right side of
\eqref{eq:W2-coupling} as the mean squared displacement of this coupling.

Put \(d=2g+1\), \(a=1-e^{-\lambda/d}\), and
\(r=e^{-\lambda/d}\).  Then
\[
 b_{g,\lambda}(\vartheta)
 =\frac{4a}{\sqrt{a^2+4r\sin^2(\vartheta/2)}}.
\]
With \(u=2\sqrt r\sin(\vartheta/2)/a\), differentiation and
\(\sup_{u\ge0}u(1+u^2)^{-3/2}=2/(3\sqrt3)<1/2\) give
\[
 |b'_{g,\lambda}(\vartheta)|\le\frac{2\sqrt r}{a}
 \le\frac{2(d+\lambda)}{\lambda};
\]
the last inequality is \(1-e^{-x}\ge x/(1+x)\).  Finally,
\cref{lem:conjugacy-quotient-metric} and \eqref{eq:W2-coupling} prove
\eqref{eq:T-Lip} on the whole group.
\end{proof}

\subsection{The transfer estimate}

For probability measures on the Polish space \((\Ptwo,\Wtwo)\), use
\[
 d_{\mathrm{BL}}(\sigma,\tau)
 =\sup_{\substack{\|H\|_\infty\le1\\ \operatorname{Lip}(H)\le1}}
 \left|\int H\,d\sigma-\int H\,d\tau\right|.
\]

\begin{proposition}[Arithmetic-to-Haar transfer]
\label{prop:arithmetic-haar-transfer}
There are an absolute constant \(C\) and, for each fixed \(\lambda>0\), a
constant \(C_\lambda\) such that, whenever \(Q>4g\),
\[
 d_{\mathrm{BL}}\bigl((T_{g,\lambda})_*\mu_{g,Q},
                       (T_{g,\lambda})_*m_g^{\mathrm{Haar}}\bigr)
 \le C_\lambda g^{-2}
   +Q^{-1/2}\exp\!\bigl(Cg^2\log(g+2)\bigr).
\tag{4.8}\label{eq:BL-transfer}
\]
Consequently the left side tends to zero whenever
\(\log Q/(g^2\log(g+2))\to\infty\).
\end{proposition}

\begin{proof}
Fix a test function \(H\) in the defining supremum and put
\(h=H\circ T_{g,\lambda}\).  Then \(h\) is central,
\(\|h\|_\infty\le1\), and by \cref{lem:T-lipschitz} it is
\(L_{g,\lambda}\)-Lipschitz with \(L_{g,\lambda}=O_\lambda(g)\).
Apply the heat bridge \eqref{eq:abstract-heat-bridge} from
\cref{app:heat-trace} with
\[
 \varepsilon_g=\frac{2(2g-1)}{\sqrt Q},\qquad
 D_g=2g^2+g,\qquad s=g^{-8}.
\]
Its smoothing term is
\[
 2L_{g,\lambda}\sqrt{2D_gs}=O_\lambda(g^{-2}).
\]
By \eqref{eq:uniform-heat-trace}, the character term is at most
\[
 Q^{-1/2}\exp\!\bigl(Cg^2\log(g+2)\bigr);
\]
all fixed numerical factors have been absorbed only after the exact heat
estimate is applied.  The resulting bound is independent of \(H\), so taking
the supremum proves \eqref{eq:BL-transfer}.  Under the stated growth
condition the logarithm of the second term is at most
\(-\tfrac12\log Q+Cg^2\log(g+2)\), which tends to \(-\infty\).
\end{proof}
